\begin{helper}[Finite exit-color selection and block partition]
Every Scheide structural configuration has a
\noderef{ScheideStructuralSelectionData}: one old block color is selected
for each exit, the selected old indices are injectively enumerated, their
complement is enumerated, and these two enumerations give the stated
unique partition of all old block indices.
\end{helper}
\begin{proof}
Put $X=\operatorname{TreeSequenceVertexSet}(T)$ and
$F=\operatorname{ScheideExitVertices}(T,\varphi)$.  The globally maximum
field of \noderef{ScheideStructuralConfiguration}, together with
\noderef{ScheideGlobalMaximumOrder}, says that $X$ is elementary in the
sense of \noderef{ColoredVertexSetElementary}.

Fix an old block index $i$.  The missing-on-base field of
\noderef{ScheideStructuralConfiguration} gives a base index $q$ such that
$C.\gamma(i)$ belongs to
\noderef{MissingColors} at $C.y(q)$.  The exact base-position field and
\noderef{TreeSequenceVertexSet} put $C.y(q)$ in $X$.  If the same color
were missing at two distinct members of $X$, those two missing-color sets
would intersect, contrary to elementariness.  Hence every old block color
is missing at exactly one vertex of $X$.

By the definition \noderef{ScheideExitVertices}, every member of $F$ is a
vertex of $X$.  For each $v\in F$, the exit-coverage field of
\noderef{ScheideStructuralConfiguration} supplies an old index $i(v)$
whose color is missing at $v$.  Make these finitely many choices.  If
$i(v)=i(w)$, then the same old block color is missing at both $v$ and $w$;
the uniqueness just proved gives $v=w$.  Thus $v\mapsto i(v)$ is
injective.  Moreover, if the selected color belonging to $w$ is missing
at $v$, its unique missing vertex is both $v$ and $w$, so $v=w$.  This is
the selected-color uniqueness clause.

Let $r=|F|$ and choose a bijection
$b:\operatorname{Fin}(r)\to\{v:v\in F\}$.  Define
$\rho(t)=i(b(t))$; the preceding paragraph proves that $\rho$ is injective
and has the required missing-color and uniqueness properties.  Let $U$ be
the complement of the range of $\rho$ in the finite type
$\operatorname{Fin}(C.s)$, put $a=|U|$, and choose a bijection
$\upsilon:\operatorname{Fin}(a)\to U$.  The range of an injection from
$\operatorname{Fin}(r)$ has $r$ elements.  It and $U$ are disjoint and
partition $\operatorname{Fin}(C.s)$, so cardinality addition gives
$C.s=a+r$.  Bijectivity of $b$ gives $r=|F|$.  Finally, membership or
nonmembership in the range of $\rho$, followed by injectivity of $\rho$
or $\upsilon$, gives the two disjoint unique alternatives for every old
index.  These data are exactly a
\noderef{ScheideStructuralSelectionData}.
\end{proof}
